\documentclass[12pt,a4paper]{article}

\usepackage[T1]{fontenc}
\usepackage[utf8]{inputenc}
\usepackage{lmodern}
\usepackage{geometry}
\usepackage{amsmath}
\usepackage{hyperref}
\usepackage{enumitem}

\geometry{margin=2.5cm}

\title{Invariantes da Estrutura de Dados Diamante Reconvergente Controlada (RDC)}
\author{Projeto RDC}
\date{\today}

\begin{document}
	\maketitle
	
	\section{Introdução}
	
	Este documento define formalmente as \textbf{invariantes testáveis} da Estrutura de
	Dados Diamante Reconvergente Controlada (RDC).
	
	As invariantes aqui descritas derivam diretamente da
	\texttt{SPEC\_RDC.md} e constituem regras
	\textbf{absolutas e não negociáveis}.
	
	Qualquer implementação que viole uma destas invariantes
	\textbf{não pode ser considerada uma implementação válida da RDC}.
	
	Este documento existe para permitir:
	
	\begin{itemize}[leftmargin=2em]
		\item Testes automatizados baseados em propriedades
		\item Auditoria arquitetural objetiva
		\item Garantia de determinismo e previsibilidade
	\end{itemize}
	
	\section{Invariantes Estruturais}
	
	\subsection{Invariante I1 — Core Único}
	
	\textbf{Definição:}  
	Existe exatamente um e apenas um Core em qualquer instância RDC.
	
	\textbf{Forma Testável:}
	\begin{itemize}
		\item A contagem de nós do tipo Core deve ser igual a 1.
	\end{itemize}
	
	\subsection{Invariante I2 — Imutabilidade do Core}
	
	\textbf{Definição:}  
	O estado do Core não pode ser modificado durante um ciclo de execução.
	
	\textbf{Forma Testável:}
	\begin{itemize}
		\item O hash ou versão do Core deve permanecer constante entre \texttt{split} e \texttt{merge}.
	\end{itemize}
	
	\subsection{Invariante I3 — Expansão Deriva do Core}
	
	\textbf{Definição:}  
	Todo nó inteligente deve ser alcançável a partir do Core.
	
	\textbf{Forma Testável:}
	\begin{itemize}
		\item Todos os nós pertencem ao fecho transitivo do Core.
	\end{itemize}
	
	\subsection{Invariante I4 — Ausência de Ciclos}
	
	\textbf{Definição:}  
	A topologia interna da RDC é acíclica.
	
	\textbf{Forma Testável:}
	\begin{itemize}
		\item O grafo interno da RDC deve ser um DAG.
	\end{itemize}
	
	\subsection{Invariante I5 — Reconvergência Obrigatória}
	
	\textbf{Definição:}  
	Todo caminho iniciado no Core deve terminar no Merge.
	
	\textbf{Forma Testável:}
	\begin{itemize}
		\item Não existem nós finais fora do Merge.
	\end{itemize}
	
	\section{Invariantes Quantitativas}
	
	\subsection{Invariante I6 — Limite Estrutural}
	
	\textbf{Definição:}  
	O número máximo de nós inteligentes ativos em uma RDC é 4096.
	
	\textbf{Forma Testável:}
	\[
	|Nodes_{\text{ativos}}| \leq 4096
	\]
	
	\section{Invariantes de Estado}
	
	\subsection{Invariante I7 — Estados Válidos}
	
	\textbf{Definição:}  
	Um Nodo-Transistor só pode assumir estados definidos pela SPEC.
	
	Estados permitidos:
	\begin{itemize}
		\item OFF
		\item STANDBY
		\item ON
		\item BOOST
		\item COLLAPSE
		\item SLEEP
	\end{itemize}
	
	\subsection{Invariante I8 — Transições Controladas}
	
	\textbf{Definição:}  
	Mudanças de estado devem ocorrer exclusivamente via gate matemático
	ou controle de energia.
	
	\section{Invariantes Matemáticas}
	
	\subsection{Invariante I9 — Dimensão do Vetor}
	
	\textbf{Definição:}  
	Todo vetor de entrada ou peso possui dimensão fixa igual a 16.
	
	\subsection{Invariante I10 — Determinismo do Gate}
	
	\textbf{Definição:}  
	Dado o mesmo vetor de entrada, o gate matemático deve produzir o mesmo resultado.
	
	\section{Invariantes de Reconvergência}
	
	\subsection{Invariante I11 — Determinismo Global}
	
	\textbf{Definição:}  
	A reconvergência deve ser determinística e independente da ordem
	de avaliação dos nós.
	
	\section{Conclusão}
	
	Estas invariantes definem o \textbf{contrato formal} da RDC.
	Qualquer violação indica erro arquitetural grave.
	Testes automatizados devem ser escritos explicitamente para cada
	invariante listada neste documento.
	
\end{document}
